% Generated draft: requires mathematical and source review.
\documentclass[11pt]{article}
\usepackage[margin=25mm]{geometry}
\usepackage{fontspec}
\setmainfont{Latin Modern Roman}
\usepackage{amsmath,amssymb,mathtools,unicode-math}
\setmathfont{Latin Modern Math}
\usepackage{enumitem,needspace,xurl,hyperref}
\setlength{\parindent}{0pt}
\setlength{\parskip}{5pt}
\setlength{\emergencystretch}{3em}
\newcommand{\N}{\mathbb{N}}
\newcommand{\Z}{\mathbb{Z}}
\newcommand{\Q}{\mathbb{Q}}
\newcommand{\R}{\mathbb{R}}
\newcommand{\C}{\mathbb{C}}
\newcommand{\F}{\mathbb{F}}
\newcommand{\A}{\mathbb{A}}
\newcommand{\PP}{\mathbb{P}}
\newcommand{\E}{\mathbb{E}}
\newcommand{\eps}{\varepsilon}
\newcommand{\dd}{\,\mathrm{d}}
\begin{document}
\section*{UnsolvedMath ID: 30004781}
\textbf{Source problem code:} catalogue entry
\subsection{Definitions and mathematical statement}
Let \(F\) be a non-Archimedean local field and let \(G\) be a connected reductive algebraic group over \(F\). The record imposes no restriction on the characteristic of \(F\), and does not assume that \(G\) is split or quasi-split.

Under the usual complex smooth representation convention, \(\Pi(G/F)\) denotes the set of isomorphism classes of irreducible admissible smooth representations of \(G(F)\) on complex vector spaces. A representation is smooth if every vector has an open stabilizer, and admissible if its subspace of vectors fixed by any compact open subgroup is finite-dimensional. The supplied record does not explicitly specify this coefficient or representation convention.

The notation \(\Phi(G/F)\) is intended to denote a set of equivalence classes of local Langlands parameters for \(G\). Its precise definition is not supplied: in particular, the parameter model, admissibility conditions, equivalence relation, and relevance condition for the particular \(F\)-group \(G\) are unspecified. Consequently, the following is a formulation schema whose target still requires clarification.

\textbf{Requested correspondence.} For every such pair \((G,F)\), construct a canonical map
\[
L_{G/F}\colon \Pi(G/F)\longrightarrow\Phi(G/F).
\]
For \(\phi\in\Phi(G/F)\), define its \(L\)-packet to be the fiber
\[
\Pi_\phi(G/F)
:=\{[\pi]\in\Pi(G/F):L_{G/F}([\pi])=\phi\}.
\]
The requested finiteness and nonemptiness condition is
\[
\forall\phi\in\Phi(G/F),\qquad
1\leq\#\Pi_\phi(G/F)<\infty.
\]
Thus the requested map is surjective with finite fibers.

\textbf{Unresolved qualifications.} For a fixed group \(G\), the nonemptiness requirement depends on the intended relevance restriction on parameters and on whether representations are taken only on \(G(F)\) or across specified inner forms. The source writes \(\Pi(G/F)\) for a single group and supplies no inner-form convention; these alternatives cannot be silently interchanged. Likewise, the word ``canonical'' is not specified by normalization, compatibility, or uniqueness requirements. The displayed fiber condition alone does not define what makes a map a local Langlands correspondence. These omissions prevent a fully determined formulation of the intended conjecture.
\subsection{Short English statement}
For every connected reductive group over a non-Archimedean local field, the source asks for a canonical assignment of representations to Langlands parameters, with each parameter corresponding to a finite, nonempty packet. The parameter conventions, relevance restrictions, and meaning of canonical remain unspecified and require human mathematical review.
\subsection{Sources}
\begin{itemize}
\item Local UnsolvedMath registry, permanent ID \(30004781\), title: \emph{The Local Langlands Correspondence for Reductive Groups}; exact source problem code: \texttt{catalogue entry}. The supplied record gives proposed year 2021 and names no proposer.
\item Supplied report citation: \emph{Automorphic Forms, Geometry and Arithmetic (2022)}, Oberwolfach Reports, pp.~2089--2156, DOI \texttt{10.4171/owr/2021/39}. \url{https://doi.org/10.4171/owr/2021/39}. The literature-source label is ``Oberwolfach Report 39/2021.'' The supplied title year 2022 and report designation 39/2021 are retained without resolving their relationship.
\item Supplied literature source: \emph{The local Langlands conjecture}. \url{https://arxiv.org/abs/2510.00632}. No author or specific supporting passage is supplied.
\item Literature assessment dated exactly \texttt{2026-08-22}: the supplied curation labels the problem ``partially solved'' and reports that the conjecture is proved for \(\mathrm{GL}_n\), with major established versions for \(p\)-adic classical groups and other families, while the stated general canonical finite nonempty parametrization remains unknown. This is supplied dated curation, not an independently verified literature assessment; the record explicitly requests independent verification of time-sensitive status and cited claims.
\item The registry labels the statement \texttt{corrected\_verified} and reports that its standalone formulation and hypotheses were checked against the official Oberwolfach report. This is a supplied source-verification claim, not verification performed for this draft.
\end{itemize}
\end{document}
